\documentclass[../../../include/open-logic-section]{subfiles}

\begin{document}

\olfileid{sth}{z}{union}
\olsection{الاتحاد}

منحنا~\olref[sth][z][sep]{prop:intersectionsexist} التقاطعات. لكن إذا
أردنا وجود الاتحادات الاعتباطية، وجب أن نضع مسلّمة أخرى:

\begin{axiom}[الاتحاد] لكل مجموعة~$A$، توجد المجموعة
$\bigcup A = \Setabs{x}{(\exists b \in A) x \in b}$.
\[
	\forall A \exists U \forall x(x \in U \liff (\exists b \in A)x \in b)
\]
\end{axiom}

يبرر التصور التراكمي-التكراري هذه المسلّمة أيضًا. لتكن~$A$ مجموعة،
فتكون~$A$ قد تكونت في مرحلة ما~$S$ (بموجب~\stageshier). وتكوّن كل
عنصر من عناصر~$A$ \emph{قبل}~$S$ (بموجب~\stagesacc)؛ ولذلك، باستدلال
مماثل، تكوّن كل عنصر من عناصر كل عنصر من~$A$ قبل~$S$. ومن ثم تكون
جميع \emph{تلك} المجموعات متاحة قبل~$S$ لكي تُكوّن في مجموعة في~$S$.
وليست تلك المجموعة إلا~$\bigcup A$.

\end{document}
